% =====================================================================
%  05-recursos.tex - Verificacion de recursos procesales
% =====================================================================

\section{Verificación de recursos procesales}\label{sec:pruebas-recursos}

Esta sección reúne la verificación del registro de recursos declarados, sus asociaciones con actividades existentes, las audiencias propias de recursos y la creación contextual de plazos.

\subsection{Pruebas focales de recursos declarados}\label{sec:pruebas-recursos-declarados}

El corte local del 19 de septiembre de 2026 verifica por separado dominio, aplicación, transporte, almacenamiento e interfaz del registro descrito en la \cref{sec:impl-recursos-declarados}. El informe \rutarepo{docs/verification-report.md} conserva la procedencia y el alcance de las ejecuciones. Las pruebas focales se complementan con aceptación HTTP y navegador contra servicios reales. Sus ejecuciones se conservan por separado; no constituyen una regresión global de una única cabeza publicada.

\begin{table}[H]
  \centering
  \caption{Evidencia focal ejecutada del registro de recursos declarados.}
  \label{tab:pruebas-recursos-declarados}
  \begin{tabularx}{\linewidth}{@{}L{2.9cm} L{2.7cm} X@{}}
    \toprule
    \textbf{Frontera} & \textbf{Resultado} & \textbf{Alcance observado} \\
    \midrule
    Dominio & 14 aprobadas & Identidades, valores, codificaciones canónicas y compromisos del recurso y del acto. \\
    Aplicación & 13 aprobadas & Permisos, reautenticación, fuentes exactas, recibos, actos y estados organizativos mediante puertos. \\
    HTTP Rust & 9 aprobadas & Rutas, cuerpos, capturas y rechazos con puertos controlados. \\
    Cliente JavaScript & 18 aprobadas & Validación de comandos, respuestas y recibos antes de consumirlos en la interfaz. \\
    PostgreSQL desechable & 12 aprobadas & Almacenamiento, atomicidad, revocación, inventario, catálogo, raíces de actos y restauración con continuación de historia. \\
    Navegador controlado & 5 y 7 casos aprobados & Dos conjuntos funcionales distintos; las dos repeticiones para capturas visuales no se cuentan como casos nuevos. \\
    Compilación web & Aprobada & Construcción de la interfaz; no sustituye aceptación integrada ni pruebas de backend. \\
    \bottomrule
  \end{tabularx}
\end{table}

\subsubsection{Identidad, capturas y confirmación}

Las pruebas de dominio comprueban límites técnicos, declaraciones explícitas y conservación de valores canónicos. Las de aplicación utilizan material verificado y lo alteran para distinguir un recibo válido de una respuesta que cambia identidad, revisión o captura. Comprueban las referencias históricas de resolución, los soportes exactos, los actos con revisión propia y la diferencia entre archivo organizativo y desistimiento declarado.

El servicio se prueba con cuatro roles y reautenticación del principal completo. Los casos de confirmación conservan el digest revisado y el recibo devuelto; una repetición exacta no requiere otro commit. Los cambios de actor, correo o rol y las respuestas incongruentes del puerto se rechazan. La comparación del contexto observado evita que el servicio sustituya una fuente seleccionada por su cabeza sin revisión humana.

La evidencia HTTP comprueba el enlace entre ruta e intención y los rechazos antes de invocar el caso de uso. Las pruebas del cliente comprueban comandos, capturas y recibos, e impiden incorporar respuestas cuya identidad o estructura no corresponde a la operación solicitada.

\subsubsection{Persistencia y límites de recuperación}

Las doce pruebas PostgreSQL aprobaron en 48.921 segundos con una base desechable y las 2328 fuentes observadas sin cambios durante la ejecución. El conjunto comprende tres casos de almacenamiento, tres de atomicidad y revocación, uno de inventario, uno de restauración, dos de catálogo y dos de raíces de actos. Son pruebas del adaptador real, distintas de los dobles de aplicación y HTTP.

Los escenarios comprueban registro, actos e historia, repetición de operación y autorización vigente. Un fallo de auditoría revierte la escritura, y una revocación durante la admisión documental impide confirmar. El inventario rechaza la sustitución de una proyección capturada; el catálogo rechaza alteraciones de columnas, restricciones, índices, guardas y privilegios. Las dos pruebas de raíz comprueban que una corrección no se acepte si, después de abrir el adaptador, su raíz se cambia para apuntar a otro acto o a una corrección en lugar de su captura inicial.

La prueba de recuperación ejecuta \texttt{pg\_dump} y \texttt{pg\_restore}, compara las capturas y conserva el correo del autor original aunque su cuenta haya cambiado. Después de restaurar, un Owner autorizado añade una revisión que continúa la historia anterior. Este resultado acredita la recuperación ensayada del adaptador. La aceptación integrada posterior con API se documenta por separado a continuación.

\subsubsection{Interfaz con transporte controlado}

El primer conjunto de Playwright aprobó cinco recorridos en 16.1 segundos del navegador. Incluye seleccionar resolución R1 aunque exista cabeza R2, registrar un acto declarado oral, archivar y reactivar a 1440 y 390 píxeles, consultar historia exacta y conciliar una respuesta incierta sin duplicar el envío. El resultado verifica interacción y contrato con respuestas HTTP controladas.

El segundo conjunto añadió siete casos de roles, revocación, cierre y conflicto. Paralegal consulta historia sin controles de gestión; Cliente no recibe navegación de recursos ni provoca consultas a esa API mediante una ruta protegida. Al revocar la asignación, la actualización retira filas y detalle. Un expediente cerrado conserva la consulta exacta sin ofrecer mutaciones habilitadas. Un conflicto conserva el borrador hasta comparar y aceptar la nueva base, y sólo entonces permite confirmar sobre la revisión esperada actualizada. La ejecución incluyó dos repeticiones adicionales destinadas a capturas: sus nueve recorridos duraron 19.2 segundos, sin convertir esas repeticiones en casos funcionales distintos.

La construcción web aprobó en 3.655 segundos. Las verificaciones locales se ejecutaron en serie, con un hilo de pruebas Rust y un trabajador de navegador. Los resultados se registran por frontera; no constituyen una nueva regresión global.

\subsubsection{Aceptación con servicios reales}

La demostración HTTP integrada aprobó en 383.388 segundos, con 2329 fuentes estables, PostgreSQL y Redis desechables, qpdf y TSA local. Registró revocación y apelación, actos orales y escritos, correcciones, archivo y reactivación. Verificó cuatro roles, aislamiento, cierre y reapertura, conflictos y repetición exacta. Antes del respaldo conservó veinte respuestas completas; después de restaurar comparó sus fuentes, autores y recibos, junto con la historia y auditoría persistidas. No se envió correo externo.

El navegador real aprobó tres recorridos en 25.3 segundos de Playwright y 264.627 segundos totales de preparación y ejecución, con 2329 fuentes estables. Owner a 1440 píxeles y Litigante a 390 registraron y corrigieron un acto oral y consultaron intactas las revisiones anteriores. Paralegal consultó sin controles de escritura, Cliente recibió denegación en la API y la revocación de membresía retiró el detalle visible. Las cuatro capturas actuales e históricas se inspeccionaron sin desbordamiento ni solapamientos. Esta evidencia verifica el prototipo con datos de ensayo; no acredita recepción judicial de un recurso.

El cierre de integración continua y la integración del registro permanecían pendientes en este corte local. La integración en la rama principal se produjo el mismo 19 de septiembre de 2026, con los controles remotos aprobados; el CI de cierre registró 2936 pruebas Rust aprobadas, ninguna fallida y una TSA externa ignorada, según \rutarepo{docs/verification-report.md}. No se midió usabilidad con personas ni cobertura global nueva. Estas pruebas tampoco verifican asociaciones de recursos con audiencias o plazos, audiencias propias ni alertas, que se comprueban por separado (\cref{sec:pruebas-asociaciones-recursos,sec:pruebas-audiencias-recursos,sec:pruebas-alertas-audiencias-recursos}), ni determinan interpretación jurídica. La matriz de la \cref{sec:matriz-recursos-declarados} conserva esa separación.


\subsection{Verificación focal de asociaciones de recursos}\label{sec:pruebas-asociaciones-recursos}

El corte local del 19 de septiembre de 2026 comprueba asociaciones a actividades existentes por fronteras independientes. Las ejecuciones fueron seriales y no constituyen una nueva regresión global. El informe versionado \rutarepo{docs/verification-report.md} conserva los resultados y las comprobaciones pendientes.

\subsubsection{Referencias, permisos y recibos}

Cuatro pruebas de dominio verificaron identidades y revisiones independientes, codificación de referencias exactas y distinción entre familia de actividad y acto opcional. La aplicación aprobó 18 pruebas: seis de flujo, ocho de límites y cuatro de tiempo. Cubren selección histórica, cabeza esperada, permisos, reautenticación, repetición de operación, cierre y conservación de fuentes e historia.

Diez pruebas HTTP aprobaron correspondencia entre URL y comando, cuerpos estrictos, límites, errores y proyecciones. El detalle de plazo conserva por separado una fecha calculada histórica y la ausencia de vencimiento operativo de la cabeza retirada. El caso adicional de acto no nulo conserva recurso R2, acto contenido en recurso R3 y soporte exacto, y rechaza una captura discordante; aprobó en 27.228 segundos con 2426 fuentes estables. Doce pruebas del cliente verificaron comandos, fuentes, recibos, alcance e invalidación de consultas pendientes.

\subsubsection{Transacción, restauración e instante de lectura}

Once pruebas PostgreSQL aprobaron almacenamiento, autorización, atomicidad, inventario, catálogo y restauración. Los rechazos comprueban aislamiento, revocación, corrupción y ausencia de escritura parcial; los positivos conservan las capturas exactas y su autoría. La restauración focal usa \texttt{pg\_dump} y \texttt{pg\_restore}, compara historia y auditoría y permite registrar otra asociación desde un Owner autorizado. No modifica las actividades vinculadas.

Una prueba adicional identificó que el instante tomado antes de esperar el bloqueo podía preceder una cabeza confirmada durante la espera. La corrección captura el instante bajo el bloqueo y conserva un corte único para detalle, página y proyección operativa. La prueba PostgreSQL específica aprobó en 15.186 segundos. Las cuatro pruebas de aplicación aceptan un corte posterior al inicio y rechazan observaciones fuera de la ventana o diferentes entre filas. Esas pruebas y las otras 14 de aplicación aprobaron juntas en 6.737 segundos, con 2426 fuentes estables.

Los doce casos PostgreSQL distintos corresponden a once positivos de la primera campaña y al positivo posterior del reloj. La primera campaña incluyó el rechazo esperado de esa prueba nueva; no se presenta como una ejecución completamente aprobada ni se suma otra vez al contabilizar el arreglo.

\subsubsection{Interacción controlada y alcance pendiente}

Doce casos distintos de navegador con HTTP controlado aprobaron el flujo de vinculación y desvinculación, selección histórica, acto opcional, separación de la actividad actual, permisos, revocación, cierre, conflicto y conciliación de una escritura confirmada cuya respuesta se perdió. Los dos casos adicionales comprueban reenvío explícito del mismo comando tras un resultado incierto no confirmado y recuperación del selector después de un conflicto al confirmar. Dos repeticiones visuales de 12.326 segundos no añaden casos funcionales distintos. Se inspeccionaron cuatro capturas de las secciones nuevas a 1440 y 390 píxeles sin desbordamiento; no se midió usabilidad con personas.

\subsubsection{Aceptación API integrada con restauración}

La campaña completa de \rutarepo{scripts/api-demo.sh} aprobó en 434.402 segundos, con 2426 fuentes sin cambios. Verificó fuentes históricas exactas frente a cabezas actuales separadas, acto opcional, desvinculación, repetición exacta, roles, revocación, cierre, archivo y aislamiento entre expedientes.

Tras restaurar se compararon 26 respuestas HTTP y se conservaron idénticas las tres raíces y cuatro revisiones de asociación del ensayo. Los instantes de lectura se validaron aparte: una consulta posterior al respaldo tiene otro instante, pero debe conservar el corte común y su vínculo con la proyección operativa. Esa diferencia no modifica las capturas históricas comparadas.

Tres escenarios distintos de navegador real aprobaron: Owner en escritorio, Litigante en móvil y permisos con Paralegal y Cliente. La primera campaña aprobó los dos flujos de escritura y detectó un error en la navegación del guion de Cliente; la repetición del recorrido corregido aprobó en 274.304 segundos, incluidos los preparativos. No se alteraron las actividades ni se duplicaron sus avisos. Los escenarios forman parte de la regresión integrada; sus cortes posteriores se registran en el informe de verificación. Las campañas de API y navegador conservan resultados separados y no equivalen a una regresión global. Tampoco acreditan el corpus jurídico (\cref{sec:pruebas-corpus-juridico}), la creación de audiencias propias de recursos (\cref{sec:pruebas-audiencias-recursos}), la activación automática, la entrega externa de avisos ni la usabilidad con personas. La matriz del \cref{sec:matriz-asociaciones-recursos} distingue esas fronteras.

\subsubsection{Consulta inversa desde actividades y alertas}\label{sec:pruebas-recursos-desde-actividad}

La consulta inversa de la \cref{sec:impl-recursos-desde-actividad} se verificó el 27 de septiembre de 2026, de forma independiente de las campañas de asociaciones del 19 de septiembre. Cinco pruebas de aplicación y seis casos distintos de PostgreSQL aprobaron su contrato; también aprobaron doce regresiones previas de persistencia. Seis pruebas HTTP nuevas y diez regresiones previas verificaron ambas familias de destino, consultas estrictas, permisos, páginas vacías y continuación. Se comprueba la separación entre capturas y cabezas actuales y se rechazan respuestas con ámbito, referencias, orden, cursor u observación incoherentes. Las seis pruebas nuevas se observaron fallar por ausencia de las rutas antes de implementar el comportamiento.

Ocho pruebas Node del cliente aprobaron validación de la página, estado, identidad del destino, capturas y descarte de respuestas posteriores a cerrar su contexto. Nueve casos distintos de navegador con HTTP controlado aprobaron la consulta y navegación; sus repeticiones focales no se contabilizan como casos nuevos. Distinguen la revisión conservada por una alerta de la revisión vinculada por la asociación, apertura histórica, regreso a actividad y bandeja, filtros y ausencia de escrituras que marquen el aviso como leído. Los rechazos y las respuestas tardías mantienen inaccesible el contenido privado. Nueve comprobaciones del planificador de fixtures reales también aprobaron.

Un recorrido de navegador con servicios reales aprobó en 11.6~s de escenario y 15.5~s del ejecutor. Se inspeccionaron sus capturas de escritorio a 1440 píxeles y móvil a 390 píxeles. Esta aceptación de navegación es independiente de la campaña API con respaldo y restauración, que también aprobó con salida cero. Las consultas inversas conservaron capturas, filtros, paginación y rechazo tras revocar pertenencia; sus instantes de lectura se validaron antes de normalizar únicamente ese dato transitorio para comparar el respaldo. Esta evidencia tampoco acredita usabilidad con personas.

Estos resultados focales y de navegador no constituyen una nueva regresión global. En ese corte permanecían pendientes el cierre global y la integración; la integración posterior, con CI, Web y Documents aprobados tanto en la revisión integrada como en su confirmación en la rama principal, se registra en \rutarepo{docs/verification-report.md}, separada de la admisión documental y de la creación contextual de plazos, verificadas también el 27 de septiembre de 2026. La revisión del manuscrito se registra de forma independiente del software. La consulta inversa no modifica el acceso documental del Cliente, cuya política conserva su decisión pendiente, ni introduce activación automática de términos; el corpus jurídico delimitado se verifica en la \cref{sec:pruebas-corpus-juridico}.


\subsection{Comprobación de la preparación y confirmación de audiencias de recursos}\label{sec:pruebas-audiencias-recursos}

La comprobación distingue valores de programación, protocolo de confirmación y persistencia transaccional. Las pruebas de dominio recorren el catálogo cerrado, su compatibilidad con recurso y modalidad escrita, la conservación de revisiones de participantes y los límites de selección. Se comprueban la admisión de 32 fichas, el rechazo de una selección superior y la duplicación de identidad aun con revisiones distintas. Un vector construido independientemente fija \texttt{RHEAR1}; las variaciones de tipo, tiempo, desfase, modalidad, lugar, nota, participantes y soporte deben modificar los bytes correspondientes. La regresión de audiencias ordinarias conserva su catálogo y formato separados.

Las pruebas de aplicación utilizan puertos controlados. Incluyen preparación sin etapa ordinaria, incompatibilidad declarada, recurso archivado, revisión de cabeza obsoleta, sustitución de captura histórica, expediente cerrado y discrepancia del soporte exacto. Los casos de autorización admiten Owner y Litigator y rechazan Paralegal y Client antes de invocar el almacén. Se contrastan selecciones de participantes ordenadas, ajenas al expediente, archivadas, ausentes o con evidencia alterada. Un cambio de principal durante la preparación impide entregar el resultado previo.

Los casos de confirmación exigen una huella idéntica antes de llamar al puerto de escritura y comprueban que éste recibe la revisión preparada completa. Una captura devuelta con evidencia válida, pero correspondiente a otra intención, se rechaza. Se alteran contenido, marcador de origen y tiempo de registro para comprobar la vinculación de la captura; el reloj debe ser UTC y no anterior a las fuentes conservadas. Una regresión intercambia procedencias y posiciones de participantes para comprobar que cada bloque queda vinculado a su identidad y revisión. La autoridad se comprueba de nuevo antes de confirmar y antes de entregar el resultado, sin equiparar un error posterior a la escritura con prueba de que ésta no ocurrió.

Una simulación en memoria conserva la creación y devuelve un error en lugar de la respuesta de confirmación. Una segunda instancia del servicio concilia expresamente la misma operación y recupera la captura original sin aumentar las escrituras. Otros casos rechazan operación o huella distintas, captura futura y pérdida de autoridad; un cambio de correo anterior a la conciliación no sustituye al autor histórico. Esta evidencia comprueba el comportamiento del llamador y del servicio, sin probar reinicio de procesos ni almacenamiento persistente. El corte anterior aprobó diez casos de dominio y 22 de aplicación. La ampliación de asociaciones aprobó seis casos de su representación de dominio, 28 del servicio de audiencias de recursos y cinco de consulta inversa. Se conservan por separado los resultados históricos de las audiencias ordinarias. Los tiempos y comprobaciones ejecutadas se registran en \rutarepo{docs/verification-report.md}. Los recorridos SQL se comprobaron en PostgreSQL 16.15 desechable, con conexiones separadas y un solo trabajador de pruebas. Nueve casos aprobaron la recuperación al reconectar después de desvincular y archivar el recurso, el rechazo de contexto cambiado, el retroceso completo ante fallo del vínculo o del origen auditable, la autorización vigente, la carrera de dos conexiones y la reconstrucción de participantes históricos. Cuatro casos adicionales aprobaron la migración repetible y el rechazo de restricciones, funciones o permisos alterados. Una regresión reprodujo antes del arreglo la pérdida simultánea de audiencia y asociación conservando su origen; la comprobación inversa del inventario y la búsqueda de operaciones previas rechazan ahora ese estado. Estos ensayos no simulan una caída eléctrica ni sustituyen una restauración completa del producto. Esas pruebas SQL no acreditaban la aceptación integrada de agenda, alertas y programación desde Qadra, documentada por separado más adelante; no se atribuyen una nueva campaña global ni resultados de despliegue. El alcance funcional y sus límites se describen en la \cref{sec:impl-audiencias-recursos}.

Las consultas se comprueban por separado en aplicación, PostgreSQL y transporte. Los casos controlados alteran ámbito, identidad, revisión, recibo, orden, continuación, reloj y principal antes de entregar datos. Los casos de PostgreSQL conservan la creación inicial después de desvincular, archivar y cerrar el expediente; también recorren revocación de pertenencia, padres o revisiones inexistentes, pérdida del origen y fallo de auditoría. El transporte contrasta el comando de confirmación y la procedencia completa de audiencia y asociación; rechaza parámetros ambiguos y respuestas que no corresponden al recurso solicitado. La prueba de composición exige que las lecturas usen la misma admisión de peticiones que el tablero. Estas comprobaciones focales no sustituyen la aceptación posterior con navegador, restauración y agenda descrita al final de este apartado.

Las pruebas de la ampliación de agenda distinguen las tres familias aun cuando comparten UUID e instante. Comprueban filtros, continuación exclusiva, conservación del origen después de desvincular, archivar y cerrar, y denegación por rol o revocación de pertenencia. También alteran la proyección persistida y simulan la pérdida del origen y el fallo de auditoría: una lectura no debe devolver una página parcial ni ocultar el error como una actividad ausente. Los resultados focales de aplicación, PostgreSQL y HTTP se registran separadamente en \rutarepo{docs/verification-report.md}; no acreditan por sí mismos el recorrido propio de navegador y restauración.

La lectura histórica desde Agenda se comprobó por separado en Qadra. Antes de implementar el cliente, dos consultas válidas de la tercera familia fueron rechazadas y la importación del cliente exacto estaba ausente. Después aprobaron 22 casos Node en 473.128 milisegundos: ocho previos de agenda y catorce nuevos. Se cotejaron los campos del resumen con la creación, su origen y asociación inicial; se probaron los extremos de cero y 32 participantes, referencias alteradas, precisión temporal de nanosegundos y descarte de respuestas tardías. La validación del transporte no recalcula los cánones criptográficos del servidor.

El navegador reprodujo la ausencia de la tarjeta propia antes de implementar el panel. La campaña posterior aprobó 17 recorridos con un trabajador en 33.3 segundos, nueve nuevos y ocho anteriores de agenda. En escritorio de 1440 píxeles y móvil de 390 se comprobó la apertura exacta de fuentes históricas y la conservación de filtros y continuación al cerrar. Otros casos verificaron el filtro propio, las tres familias con igual identidad, la pérdida de permiso y el descarte de una respuesta retenida al actualizar, cambiar el intervalo, abandonar la vista o terminar sesión. Las respuestas HTTP fueron controladas: estos resultados no sustituyen aceptación con servicios reales, restauración, cobertura global o despliegue. Estas pruebas de consulta no acreditan por sí solas la programación ni las alertas.

El cliente de programación y la ampliación de asociaciones se verificaron con desarrollo dirigido por pruebas. La primera ejecución del cliente tuvo doce fallos por operaciones ausentes y nueve lecturas aprobadas; el listado propio reprodujo cuatro fallos y las asociaciones cuatro rechazos de la familia nueva. Después aprobaron 42 casos Node en 1261.504 milisegundos, con ejecución secuencial, incluyendo las regresiones anteriores. Comprobaron alcance, fuentes, autor, cabeza, orden, cursor y conservación del comando ante respuestas inciertas. La consulta exacta no produjo otro envío y sólo confirmó material coincidente.

El helper de recuperación añadió ocho casos, inicialmente bloqueados por la importación ausente y después aprobados en 218.317 milisegundos. Se verificaron autorización antes de renovar referencias, conservación de campos incompletos e identidades, descarte de respuestas tardías y rechazo del borrador de otra identidad completa. Un resultado ausente no eliminó el envío incierto. Estas pruebas locales del cliente no sustituyen aceptación real, restauración o integración global.

La interfaz de programación añadió dieciséis recorridos de navegador con transporte controlado y un trabajador: doce aprobaron en una campaña parcial de 38.5 segundos y los cuatro restantes en 17.8 segundos después de corregir una proyección simulada de participante. La regresión de Agenda y asociaciones conservó trece casos aprobados en 32.7 segundos. Se verificaron las anchuras de 1440 y 390 píxeles, selección histórica de recurso y acto, soporte obligatorio, participante exacto, comparación explícita de cabeza, preparación antes de confirmar y conservación del comando ante un resultado incierto. La consulta de recuperación no repite el envío; su repetición requiere una decisión expresa. Se comprobaron también lectura de historia desvinculada, denegación de acceso, cierre y archivo.

Una prueba nueva reprodujo un defecto del segundo intento de recuperación de borrador: después de fallar la lectura de una fuente, el formulario podía habilitarse sin consultarla de nuevo. El caso falló al observar una sola consulta frente a dos esperadas. La corrección exige renovar las fuentes antes de desbloquear en cada intento; dos respuestas temporales de error mantienen el borrador inaccesible y una tercera respuesta válida permite recuperarlo sin escritura. Los recorridos incluyen expiración, reingreso de la misma cuenta, descarte por cambio de cuenta y salida explícita. Estas verificaciones no sustituyen la aceptación con servicios reales, restauración, alertas o integración completa.

La compilación Web aprobó en 3.25 segundos. Las capturas de revisión y resultado en escritorio y móvil conservaron legibilidad y ausencia de desbordamiento horizontal. Sólo esos dos casos positivos se repitieron para retirar el foco antes de obtener las capturas; aprobaron en 12.6 segundos y no se contaron como casos adicionales. Se mantuvo la advertencia existente sobre tamaño del paquete Web, sin modificar su umbral.

\subsubsection{Aceptación propia con servicios reales y restauración}

El recorrido HTTP propio conserva dos audiencias con autores distintos, fuentes históricas, asociación inicial y recibos estables. Verifica preparación sin escritura, confirmación y repetición exactas, conflicto de cabeza, permisos vigentes y conservación después de actualizar participantes, desvincular, archivar y cerrar. La captura focal aprobó esos controles. Posteriormente se compararon todas las tablas públicas y los valores lógicos de las secuencias antes y después de respaldar y restaurar con PostgreSQL 16.15. El procedimiento invalidó las sesiones antiguas antes de iniciar el servidor restaurado y exigió MFA nuevo; conservó los comandos repetidos, la agenda cerrada y la misma ocurrencia y lectura de alerta. La continuación final aprobó en 41.265 segundos, sin repetir la creación ni la compilación. No se trata de una simulación de caída eléctrica ni de restauración de producción.

La campaña expuso un timeout real del bloqueo global de auditoría. La planificación reconstruía la misma audiencia y sus fuentes por cada destinatario y anticipación. El muestreo identificó una transacción de 1.361 segundos que bloqueaba Agenda. Una regresión determinista reprodujo 65 verificaciones completas frente a las cinco de la pasada inicial. Reutilizar la evidencia validada sólo durante la misma transacción eliminó esa multiplicación, conservando las comprobaciones de cada plan. Los dos casos nuevos aprobaron en 14.61 segundos y tres regresiones de integridad en 66.02 segundos. Se rechazaron ocho alteraciones con checksum coherente y la pérdida del marcador entre transacciones. Cuatro consultas posteriores sobre la base conservada aprobaron entre 0.192 y 0.281 segundos. Los límites de espera y la criptografía permanecieron intactos; la medición no se extrapola al producto completo.

Con la corrección aplicada, la demostración HTTP compuesta completa aprobó en 442.355 segundos, incluida una compilación caliente de 0.69 segundos. El ensayo creó datos desechables y terminó los recorridos previos, el respaldo y la restauración en una sola ejecución. Conservó las capturas propias, autores, asociaciones, comandos repetidos, Agenda y una sola ocurrencia leída después de reiniciar el generador. Los controles existentes de roles, documentos, plazos, alertas y miembros también aprobaron; el cierre preservó cuatro documentos, 70 eventos de auditoría y el ZIP de evidencia idéntico. Se retiraron los servicios de prueba. Este resultado cierra la aceptación API compuesta, sin atribuirle validación de navegador ni despliegue.

Dos recorridos de navegador con servicios reales aprobaron en 52.6 segundos, con un trabajador, a 1440 y 390 píxeles. Comprobaron programación, pérdida de la respuesta de creación, conciliación expresa después de desvincular, lectura del mismo origen en Agenda y Alertas y recibo de lectura sin segundo envío de creación. La ejecución inicial descubrió una carrera del test: esperaba un encabezado que podía aparecer antes del POST. La prueba corregida espera su finalización conservando el resultado 201 y las demás aserciones. Las cuatro capturas visuales se inspeccionaron sin nuevos solapes ni desbordamiento horizontal. Este navegador precedió al ajuste del scheduler; la revisión final conserva pendientes sus gates globales, integración y despliegue. El registro técnico mantiene separados los fallos iniciales, las correcciones de los instrumentos y la evidencia posterior.

\subsection{Verificación de la creación contextual de plazos}\label{sec:pruebas-plazo-contextual}

La aceptación de este flujo debe comprobar que preparar no crea ninguno de los dos registros y que confirmar conserva una única operación, sus recibos y las capturas exactas seleccionadas. Incluye repetición idéntica, denegación por cambio de permisos o dependencias, rechazo de operaciones ordinarias preexistentes y concurrencia. La atomicidad exige además provocar un fallo durante la transacción y comprobar que no quedan escrituras parciales; rechazar una precondición antes de escribir no demuestra por sí solo esa propiedad.

El 27 de septiembre de 2026 la demostración HTTP integrada terminó con salida cero: comprobó confirmación conjunta, repetición exacta, consulta en agenda y rechazo de una preparación obsoleta. Después de restaurar PostgreSQL y renovar la sesión, conservó las capturas, los recibos del plazo y la asociación, y su proyección en agenda. Las comprobaciones focales del backend aprobaron cinco casos de aplicación, nueve distintos de PostgreSQL y tres de HTTP; no constituyen una regresión global. La interfaz aprobó once verificaciones Node, once recorridos con HTTP controlado en 28.8 s y uno con servicios reales en 13.6 s de escenario (17.8 s del ejecutor de navegador). Se inspeccionaron las capturas de escritorio y móvil, sin desbordamiento horizontal. La regresión global permanece pendiente. La compilación y revisión del PDF se registran por corte en el informe de verificación. La evidencia de la \cref{sec:impl-plazo-contextual} se registra por separado del tablero en \rutarepo{docs/verification-report.md}.
